\section{Track Geometry Model}
\label{sec:track-geometry}

The cornering model of \cref{sec:cornering-model} requires a curvature profile
$\kappa(s)$ describing how sharply the track turns at each point along its
length. This profile is generated by \texttt{TrackShape} from a small number of
geometric parameters and is also used purely for visualizing the track shape in
the application's user interface.

\subsection{Idealized track shape}
\label{sec:idealized-shape}

The track is modeled as an idealized oval: two straight sections and two
semicircular corners, in the style of a running track or velodrome. Given total
track length $L$ (\texttt{track\_length}) and a dimensionless fraction
$c \in [0.20, 1.0]$ (\texttt{corners}) specifying what portion of $L$ is
devoted to cornering, the arc length of each straight and each corner is
\begin{equation}
    \ell_{\mathrm{straight}} = \frac{L(1-c)}{2}, \qquad
    \ell_{\mathrm{corner}} = \frac{Lc}{2}.
    \label{eq:straight-corner-length}
\end{equation}
The user-facing control for $c$ is labeled from \enquote{hotdog} at its minimum
($c=0.20$, long straights and tight, brief corners) to \enquote{circle} at its
maximum ($c=1.0$, no straights at all, i.e.\ a single circular track).

\subsection{Corner radius}
\label{sec:corner-radius}

Each corner is modeled as an exact semicircular arc: a $180^\circ$ turn subtending
arc length $\ell_{\mathrm{corner}}$. Since arc length equals radius times swept
angle, and the swept angle of a semicircle is $\pi$,
\begin{equation}
    r = \frac{\ell_{\mathrm{corner}}}{\pi}.
    \label{eq:corner-radius}
\end{equation}
The corner curvature is then simply $\kappa_{\mathrm{corner}} = 1/r$.

\subsection{Curvature profile and transition smoothing}
\label{sec:curvature-smoothing}

Without smoothing, curvature along the track would be a step function: zero on
the straights, $1/r$ on the corner arcs, with a discontinuous jump at each
straight-to-corner boundary. Physically, no rider leans instantaneously; the
transition into and out of a corner takes some finite distance. \texttt{TrackShape}
approximates this by passing the raw step-function curvature through a
box (moving-average) filter of width \texttt{transition} (default
\SI{30}{\metre}), producing a ramp-plateau-ramp curvature profile at each corner
entry and exit rather than a sharp step. This is not a true clothoid (an Euler
spiral, the curvature-linear-in-arc-length transition used in road and rail
design) but serves an analogous smoothing purpose with substantially less
implementation complexity, at the cost of a curvature profile whose transition
shape is set by the filter kernel rather than by an exact geometric construction.
It is this smoothed transition ramp that motivates the \texttt{max\_step}
restriction discussed in \cref{sec:numerical-integration}.

This smoothing also leaves the track's total accumulated turning unchanged. A
box filter is a normalized moving average, so it redistributes curvature
locally without changing its integral, provided the signal is unchanged
wherever the averaging window reaches past the domain's ends -- true here,
since both track endpoints ($s=0$ and $s=L$) fall at the midpoint of a
straight, well away from any corner, where the raw curvature is already zero.
Consequently
\begin{equation}
    \int_0^L \kappa(s)\,ds = 2\pi
    \label{eq:total-turning}
\end{equation}
exactly, independent of the transition length: each of the two corners is an
exact semicircle contributing $\kappa_{\mathrm{corner}}\,\ell_{\mathrm{corner}}
= \pi$ radians of turning by \cref{eq:corner-radius}, and smoothing only moves
\emph{where} that turning happens along $s$, not how much of it there is. This
holds for any track length $L$ and corner fraction $c$, and is exactly why the
reconstructed centerline of \cref{sec:position-reconstruction} always closes
into a loop that returns to the starting point after one full lap, whatever
transition length is configured.

\subsection{Position reconstruction}
\label{sec:position-reconstruction}

For visualization only -- this reconstruction has no effect on the physics of
\cref{sec:simulation-model,sec:cornering-model} -- the track's heading and
Cartesian centerline coordinates are recovered from the curvature profile by the
standard Frenet-style construction: heading is the running integral of
curvature, and position is the running integral of the heading's cosine and
sine components,
\begin{align}
    \phi(s) &= \int_0^s \kappa(s')\,ds', \label{eq:heading} \\
    X(s) &= \int_0^s \cos\phi(s')\,ds', \label{eq:x-pos} \\
    Y(s) &= \int_0^s \sin\phi(s')\,ds'. \label{eq:y-pos}
\end{align}
Each integral is evaluated numerically by cumulative trapezoidal integration.
As a consequence of \cref{eq:total-turning}, the reconstructed centerline
closes on itself: after one full lap, $(X,Y)$ returns to within \SI{1}{\metre}
of the starting point, for any track length or transition setting.
